\chapter{Conventions and Units}
\label{ch:conventions_units}

This chapter collects the unit systems, notation, and Fourier
conventions used in this document.  The harmonic force
constants may originate either from phonopy finite-displacement
calculations or from Quantum ESPRESSO (QE) density-functional
perturbation theory, while the QuaTrEx solver uses its own internal
frequency-squared convention.  Since the later extraction, transport,
and anharmonic self-energy steps all rely on consistent units and phase
conventions, we summarize them here once and refer back to this section
% REVIEW: this file IS the chapter; refer to the document, not the chapter
throughout the document.

% ----------------------------------------------------------------------
\section{Units in phonopy and Quantum ESPRESSO}
\label{sec:phonopy_units}
% ----------------------------------------------------------------------

Phonopy stores force constants in eV/\AA$^2$, atomic masses in amu, and
atomic positions in \AA.  The dynamical matrix eigenvalues therefore
have units
\begin{equation}
  [\lambda] = \frac{\mathrm{eV}}{\mathrm{amu}\cdot\mathrm{\si{\angstrom}}^2}.
\end{equation}
Phonopy converts these eigenvalues to frequencies in THz via
% REVIEW: the appended ~15.633 THz is the VaspToTHz prefactor alone (lambda-independent), not the value of f
\begin{equation}\label{eq:VaspToTHz}
  f\;[\mathrm{THz}]
  = \mathrm{sgn}(\lambda)\,\sqrt{|\lambda|}
    \times
    \frac{1}{2\pi}\sqrt{\frac{e}{m_u A^2}}  \approx 15.633\;\mathrm{THz},
\end{equation}
where $e = 1.602\times 10^{-19}$\,J is the elementary charge,
$m_u = 1.661\times 10^{-27}$\,kg is the atomic mass unit, and
% REVIEW: 1e-10 m is the Angstrom-to-metre factor, not metre-to-Angstrom
$A = 10^{-10}$\,m is the \AA-to-metre conversion factor.

When phonopy is used together with the QE interface, the internal
representation inherited from QE uses atomic units: lattice vectors are
in bohr and force constants in Ry/bohr$^2$.  These quantities are
therefore converted before being used together with the rest of the
workflow.  The required relations are
\begin{align}
  a\;[\mathrm{\si{\angstrom}}]
    &= a\;[\mathrm{bohr}] \times a_0,
    & a_0 &= 0.52918\;\mathrm{\si{\angstrom}/bohr}, \\
  \Phi\;[\mathrm{eV/\si{\angstrom}^2}]
    &= \Phi\;[\mathrm{Ry/bohr^2}] \times \frac{E_h/2}{a_0^2},
    & \frac{E_h/2}{a_0^2} &\approx 48.59.
\end{align}
Here $E_h$ denotes the Hartree energy, so that $E_h/2 = 1$\,Ry.  In the
same way, Cartesian forces parsed from QE output are converted using
\begin{equation}
  1\;\mathrm{Ry/bohr} \approx 25.711\;\mathrm{eV/\si{\angstrom}}.
\end{equation}

The downstream workflow therefore treats eV, \AA, and amu as the common
intermediate unit system, independent of whether the original data came
from phonopy directly or through QE.  In other words, everything is
normalized to the phonopy eV/\AA/amu convention upon loading.

% ----------------------------------------------------------------------
\section{QuaTrEx solver units}
\label{sec:solver_units}
% ----------------------------------------------------------------------

The QuaTrEx phonon solver operates in angular frequency squared,
\begin{equation}
  \omega^2 \quad \text{in} \quad (\mathrm{rad/s})^2.
\end{equation}
Accordingly, the mass-weighted real-space blocks passed to the solver
must be converted from the phonopy eigenvalue units
eV/(amu$\cdot$\AA$^2$) to $(\mathrm{rad/s})^2$.  The corresponding
conversion factor is
\begin{equation}\label{eq:conversion}
  C = \frac{e}{m_u A^2}
  \approx 9.648 \times 10^{27}\;\mathrm{(rad/s)^2}.
\end{equation}
If $\Phi(0,\vc{R})$ denotes a harmonic force-constant block in
eV/\AA$^2$, then the real-space matrix block used by QuaTrEx is
\begin{equation}
  H(\vc{R})
  = C\,
    \frac{\Phi(0,\vc{R})}{\sqrt{m\,m'}},
\end{equation}
with $m$ and $m'$ in amu.  The resulting $H(\vc{R})$ therefore has units
of $(\mathrm{rad/s})^2$.

The same logic extends to cubic force constants.  After mass-weighting,
the third-order tensor
\begin{equation}
  \Psi_{\alpha\beta\gamma}
  = \frac{\Phi^{(3)}_{\alpha\beta\gamma}}
         {\sqrt{M_\kappa M_{\kappa'} M_{\kappa''}}}
\end{equation}
has units
\begin{equation}
  [\Psi] = \frac{\mathrm{eV}}
                 {\mathrm{\si{\angstrom}}^3\cdot\mathrm{amu}^{3/2}}.
\end{equation}
Its SI conversion factor is
\begin{equation}\label{eq:conversion_fc3}
  \texttt{CONVERSION\_FC3}
  = \frac{C}{A\sqrt{m_u}}
  = \frac{e}{m_u^{3/2} A^3}
  \approx 2.368\times10^{51},
\end{equation}
which converts
eV/(\AA$^3\cdot$amu$^{3/2}$) to the SI units required in the
phonon--phonon self-energy expressions.

% REVIEW: code->doc -- the production phonon-phonon stack (quatrex/phonon/units.py; static_self_energy.py device_fc3_mass_weighted; bubble_prefactor_thz = i*hbar*dOmega/(4 pi)) actually runs in a THz^2 / plain-THz internal system, not (rad/s)^2; reconcile
\paragraph{THz-based internal unit system.}
The $(\mathrm{rad/s})^2$ statement above is the SI-facing equivalent.  In
practice the production phonon--phonon stack
(\texttt{quatrex/phonon/units.py},
\texttt{static\_self\_energy.py:\,device\_fc3\_mass\_weighted}) works in a
\emph{THz-based} internal system: the dynamical matrix $D$ and all
self-energies are carried in $\mathrm{THz}^2$, while the frequency axis
$\omega$ is in plain THz, with the bubble assembled using
$\texttt{bubble\_prefactor\_thz} = i\hbar\,\Delta\Omega/(4\pi)$.  The
$\mathrm{THz}^2$ convention and this prefactor are documented in
\cref{ap:fa_overview}; the $(\mathrm{rad/s})^2$ form and
\cref{eq:conversion_fc3} provide the SI-facing equivalents.

\paragraph{Green's-function sign convention.}
% REVIEW: labeling tension -- theory/30_scba_eta0.tex and infrared.tex write G^<_eq=-i n A (textbook sign) while invoking "occupation-positive"; the stored solver sign is +i n A. Those two appendix statements should cite the textbook sign G^<=-inA explicitly. (FLAG-ONLY: no math changed here.)
The solver stores the phonon Green's functions and self-energies
\emph{occupation-positive}: $-iG^{\lessgtr}(\omega)\succeq0$ for $\omega>0$, the
same sign as the lead injection
$\Sigma_\ell^{\lessgtr}=+i\,n_\ell(\pm1)\Gamma_\ell$. This differs from the
common textbook choice $G^<=-i\,nA$ by an overall sign of $G^{\lessgtr}$ --- and
hence of any self-energy quadratic in $G$, such as the three-phonon bubble; see
the remark following \cref{eq:sigma_r} for the consequent sign handling.

% ----------------------------------------------------------------------
\section{Notation and Fourier conventions}
\label{sec:notation_fourier}
% ----------------------------------------------------------------------

We first collect the symbols used throughout this work.

\begin{center}
\begin{tabular}{llr}
  \toprule
  Symbol & Meaning & Range \\
  \midrule
  $l,m$                              & Supercell index                              & $\Ncell$ \\
  $\kappa,\kappa',\kappa''$          & Basis-atom within a primitive cell            & $\nprim$ \\
  $\alpha,\alpha',\alpha''$          & Cartesian component ($x,y,z$)                & $C=3$ \\
  $i, j, k$                          & DOF index                                    & $\ndof$ \\
  $\mu, \nu$                         & Primitive-cell DOF index                     & $3\nprim$ \\
  $\vq_\perp$                        & Transverse wavevector (2D BZ)                & $N_q$ points \\
  $x,y$                              & Transport-direction DOF                      & $\ntrans = 3\nprim N_\mathrm{transport}$ \\
  $\vR_\perp$                        & Transverse lattice vector                    & $N_q$ points (dual) \\
  $\omega,\omega'$                   & Real frequency                               & $N_\omega$ points \\
  $\tau$                             & Conjugate time variable                      & $N_\tau = N_\omega$ \\
  $N$                                & Total number of $\vq$-points, $N_q$          & \\
  \bottomrule
\end{tabular}
\end{center}

Here $\Ncell$ is the total number of supercells (e.g.\ $3^3 = 27$),
$\nprim$ the number of atoms in the primitive cell (e.g.\ $2$ for
bulk Si), and
% REVIEW: Ncell is the TOTAL cell count (=27), so total DOF = 3 Ncell nprim; drop the spurious cube
$\ndof = 3\nsuper = 3\Ncell\nprim$ the total number of degrees of
freedom in the simulation.  $\ntrans$ is the total number of DOFs in
the transport direction (i.e.\ when transverse cells are not
included), and $N_\mathrm{transport}$ is the number of transport
(super)-cells.

Bold symbols such as $\vR$ and $\vq$ denote lattice vectors and
wavevectors, respectively.  The displacement of atom $\kappa$ in
cell $l$ along direction $\alpha$ is written $u_\alpha(l\kappa)$.
The lattice vector of cell $l$ is $\vR_l$, while $\vc{\tau}_\kappa$
denotes the basis-atom position within the unit cell.  Unless stated
otherwise, $\vc{\tau}_\kappa$ and $\vq$ are expressed in fractional
% REVIEW: two phase conventions actually coexist; the blanket "2pi in all phases" claim contradicts eq:dft_fwd/eq:dft_inv -- state both explicitly
crystal coordinates.  Two distinct phase conventions are used in this
work: the dynamical-matrix transforms
(\cref{eq:DA_conv,eq:DB_conv,eq:fourier_series_conv}) carry an explicit
$2\pi$, i.e.\ $e^{+2\pi i\,\vq\cdot\vR}$, whereas the
Green's-function/self-energy discrete Fourier pair
(\cref{eq:dft_fwd,eq:dft_inv}) absorbs the $2\pi$ into $\vq$ and uses the
opposite forward sign, $e^{-i\,\vq\cdot\vR}$.

The harmonic force constants are written as
\begin{equation}
  \Phi_{\alpha\alpha'}(0\kappa,\,l'\kappa'),
\end{equation}
and the corresponding mass-weighted dynamical matrix at wavevector
$\vq$ is a Fourier transform of these real-space couplings.
Two phase conventions are used in this work.  In the phonopy
convention,
\begin{equation}\label{eq:DA_conv}
  D_{\alpha\alpha'}^{(\mathrm{A})}(\kappa\kappa',\vq)
  = \frac{1}{\sqrt{m_\kappa m_{\kappa'}}}
    \sum_{l'} \Phi_{\alpha\alpha'}(0\kappa,\,l'\kappa')\,
    e^{2\pi i\,\vq\cdot
      [\vR_{l'} + \vc{\tau}_{\kappa'} - \vc{\tau}_\kappa]},
\end{equation}
whereas the QuaTrEx convention uses only the lattice vector in the
phase,
\begin{equation}\label{eq:DB_conv}
  D_{\alpha\alpha'}^{(\mathrm{B})}(\kappa\kappa',\vq)
  = \frac{1}{\sqrt{m_\kappa m_{\kappa'}}}
    \sum_{l'} \Phi_{\alpha\alpha'}(0\kappa,\,l'\kappa')\,
    e^{2\pi i\,\vq\cdot\vR_{l'}}.
\end{equation}
The two are related by a diagonal unitary gauge transformation,
\begin{equation}\label{eq:gauge_conv}
  D^{(\mathrm{B})}(\vq)
  = P(\vq)\,D^{(\mathrm{A})}(\vq)\,P^\dagger(\vq),
\end{equation}
with
\begin{equation}
  P(\vq)_{\kappa\alpha,\kappa'\alpha'}
  = \delta_{\kappa\kappa'}\delta_{\alpha\alpha'}\,
    e^{2\pi i\,\vq\cdot\vc{\tau}_\kappa}.
\end{equation}
Under Convention~B, the dynamical matrix is a pure lattice Fourier
series,
\begin{equation}\label{eq:fourier_series_conv}
  D^{(\mathrm{B})}(\vq)
  = \sum_{\vR} H(\vR)\,e^{2\pi i\,\vq\cdot\vR}.
\end{equation}
% REVIEW: symbol overloading -- H(R) here is the (rad/s)^2 block of eq:conversion (= C*Phi/sqrt(mm')), i.e. it carries the C conversion, unlike the bare mass-weighted Phi/sqrt(mm') used in eq:DB_conv
Here $H(\vR)$ is the SI-converted mass-weighted block introduced in
\cref{eq:conversion}, $H(\vR) = C\,\Phi(0,\vR)/\sqrt{m\,m'}$, so that it
already carries the conversion factor $C$; the $D^{(\mathrm{B})}$ of
\cref{eq:DB_conv} instead uses the bare mass-weighted force constants
$\Phi/\sqrt{m_\kappa m_{\kappa'}}$ and therefore differs from
\cref{eq:fourier_series_conv} by the overall factor $C$.

For the transverse directions we write $\vq_\perp$ for the
two-dimensional transverse wavevector.  A partial Fourier sum over
transverse lattice vectors $\vR_\perp$ then takes the form
\begin{equation}
  D(\vq_\perp)
  = \sum_{\vR_\perp} M(\vR_\perp)\,
    e^{2\pi i\,\vq_\perp\cdot\vR_\perp}.
\end{equation}
The transverse wavevectors are sampled on a Monkhorst--Pack mesh,
\begin{equation}
  q_{\perp,i} = \frac{n_i + 0.5}{N_i} - 0.5 + \mathrm{shift}_i,
  \qquad n_i = 0,\ldots,N_i-1.
\end{equation}

We also define the discrete Fourier-transform pair
\begin{align}
  G(\omega;\vq) &= \sum_{\vR} g(\omega;\vR)\,e^{-i\vq\cdot\vR},
  \label{eq:dft_fwd} \\
  g(\omega;\vR) &= \frac{1}{N}\sum_{\vq} G(\omega;\vq)\,e^{+i\vq\cdot\vR},
  \label{eq:dft_inv}
\end{align}
with the associated discrete convolution theorem
\begin{equation}\label{eq:conv_theorem}
  \frac{1}{N}\sum_{\vq'} A(\vq')\,B(\vq-\vq')
  = \sum_{\vR} a(\vR)\,b(\vR)\,e^{-i\vq\cdot\vR}.
\end{equation}